% Link: https://www.jobs.manchester.ac.uk/Job/JobDetail?JobId=32976&Source=JobtrainGoogle 

% what the say:

% THE PROJECT

% Rules of life that govern microbialccommunities remain largely unknown, 
% limiting our ability to predict or harness them.

% Our consortium will study the microbial communities inhabiting geothermal
% springs in Iceland to discover the rules of microbial life. 

% High temperatures and extreme pH and lack any organic carbon input,
% instead using volcanic gases (CO2, H2 etc) released by volcanic activity.

% Our goal is to use understanding of the ecology, evolution and molecular biology 
% of these microbial communities to harness them for new technologies to 
% capture carbon from industrial waste gases for bioconversion into high-value chemicals.


% THE POSITION

% Purpose of this role is to use mathematical and computational modelling combined with
% experimentation to understand the fundamental eco-evolutionary processes that shape the
% composition, dynamics, and function of carbon-fixing hot-spring communities.

% The PDRA will use a combination of ecological time-series inference methods and metabolic modelling alongside
% high-throughput bioreactors to test model predictions experimentally. Through an iterative loop
% of modelling and experimentation we will build a new body of eco-evolutionary theory explaining
% the rules of microbial life

% The PDRA will build and analyse mathematical models, statistically analyse large datasets, 
% and help plan (and possibly implement) new experiments.
% They will also be responsible for the timely preparation of interim reports and will lead the
% drafting of manuscripts for publication. The PDRA will disseminate the project findings to
% scientists through oral presentations at conferences/workshops and to the public via outreach
% festivals/school visits

% PRDA: Postdoctoral research assistant

% FROM THEIR WEBSITE
% Be positive! Instead of focusing on any gaps in your experience, 
% focus on what you’ve achieved already and set out actions you’ve identified to develop further. 


% think about how to group criteria together under headings 
% and provide an example that evidences a few at once.
% describe the SITUATION, the TASK you had to complete, 
% the ACTION you took, and the RESULT. 
% (Pro tip: this is an excellent way to prepare for your interview too.)



My understanding is that this is a highly interdisciplinary position. 
With the following letter, I will try to highlight my skills and experience 
that I believe are relevant to the position, 
by addressing the person specification criteria in groups.


\textbf{Person Specification}

\textbf{Essential Skills, Knowledge and Experience}

\textit{
    Have, or be about to obtain, a PhD in a relevant field 
    (e.g. ecology, evolutionary biology, microbiology, mathematical biology etc)
}

    My PhD focuses on microbial ecology and challenges derived from 
    high throughput sequencing data representing microbiome and their analysis.

    Beyond developing workflows and software tools tailored to 
    real-world microbial communities, particularly for the analysis amplicon and shotgun
    metagenomics data, I also served as a lead developer of a \href{https://prego.hcmr.gr/}{data integration platform}
    designed to extract and consolidate associations among species, environments,
    and biological processes from both literature and abundance tables from key platforms
    such as MGnify, and IMG/JGI.

    Moreover, besides developing tools, I also applied them to real world data.
    Together with my partner in science and life, we studied a hypersaline microbial swamp in Karpathos, Greece.
    She did the sampling and the lab work, while I managed to reconstruct a set of metagenome assembled genomes (MAGs)
    and have some first insight on the mechanisms these strains used to adapt to such extreme salinity.


    My first two years side project gave me the chance to work with a large phylogenetic tree that 
    we used to distinguish non-eukaryotic COI sequences from eukaryotic ones. 

    Last, in my last two years side project, I became aware of metabolic modeling and its potential,
    especially in the study of microbial interactions; my long-standing interest of mine. 
    During this last time of my PhD, we developed a Multiphase Monte Carlo Sampling (MMCS) algorithm 
    for efficient flux sampling, in collaboration with Dr. Chalkis, 
    a computational geometer, and very good friend of mine.
    Since then, metabolic modeling has become my major tool to approach my favorite research question: 
    microbial interactions.


\textit{
    Strong conceptual understanding of eco-evolutionary theory,
    particularly as applied to microbial communities
}

    Prof. Vellend's work was a turning point in my understanding of community ecology.
    Based on his idea on all self-replicating systems being shaped by four fundamental processes,
    (selection, drift, dispersal, speciation) 
    any such system can be studied in an evolutionary framework.

    This is how I think of applying metabolic modeling and in microbial communities, too.
    Even since now, it has been widely used to study evolution at the strain level - 
    my goal is to apply its methodologies in eco-evolutionary question. 

    In a recent study of ours, we were able to show the limitations that 
    except for curated Genome Scale Models (GEMs), predicted growth rates and their ratios 
    thus, the interaction strength between the two taxa predicted, 
    do not correlate with such obtained from in vitro data~\cite{joseph2024predicting}.
    To address this limitation, we have just released the first beta-version of a microbial growth data repository, 
    called \href{https://mgrowthdb.gbiomed.kuleuven.be/}{{$\mu$}growthDB}
    The broader the endorsement by the microbiology community, 
    the more effectively it will enable us to constrain models with higher accuracy.

    On top of this, \textit{microbetag} (accepted in principle in \textit{Genome Biology}) 

    \cite{zafeiropoulos2024microbetag}

    In, I was able to submit a proposal to the 
    European Molecular Biology Organization (EMBO) postodoctoral fellowship program,
    where I proposed a project to study the relationship between
    resource preference hierarchy and the evolution of auxotrophies in microbial communities.
    Although my proposal was eventually not selected, it made it to the interviews step and 
    received very positive feedback from the reviewers.




    is a tool that annotates microbial co-occurrence networks

, for example StrategAux





\textit{Experience designing and analysing mathematical models of biological systems}

In our perspective paper "Metabolic models of human gut microbiota: Advances and challenges"
we highlight how one can a set of different modeling approaches using a toy model of three common 
gut species with different metabolic capabilities: a carbohydrate fermenter, an acetogen, and a butyrate producer.

We display how one can use kinetic models, 
inability to handle flexible metabolic responses means that in many cases it does not meet
Einstein's famous 'as simple as possible but not simpler' criterium


\textit{Experience testing models with data}

I am now working with Salmonella in vitro data to decipher why a set of probiotics decrease Salmonella. 


\textit{A keen interest in developing new eco-evolutionary theory}

My main interest lies with StrategAux: my hypothesis that auxotrophies a species adapts is shaped 
from their resource hierarchy (nutrients' preference). 
Coalescence can go my synthesis a step further, since invasion would add dynamics in my hypothesis. 



% --------------------------------------------
\newpage\vspace*{0.3in} 
% --------------------------------------------

% overall productivity and the impact of your publications.
\textit{Strong journal publication record.}
% evidence that you consistently publish in peer-reviewed venues.
\textit{Track record of peer-reviewed scientific publication(s)}

To date, I have co-authored 17 peer-reviewed publications, 
serving as first author on 9 and corresponding author on 4.

My PhD related work, was mostly published in the GigaScience journal,
a well respected journal in the field of bioinformatics and computational biology.
Both pema~\cite{zafeiropoulos2020pema} and metagoflow~\cite{zafeiropoulos2023metagoflow}
as well as the review I wrote on the IMBBC High Performance Computing (HPC) system 
~\cite{zafeiropoulos20210s}.

My two side projects at that time were also published in peer-reviewed journals:
the COI classifier~\cite{zafeiropoulos2021dark}, where we were able to show that 
non-eukaryotic sequences can be sequenced when sequencing for COI, was published in the Metabarcoding and Metagenomics journal,
a rather new journal at the time. 
While our MMCS algorithm for efficient flux sampling, made it to the top conference in computational biology,
Symposium on Computational Geometry~\cite{socg2021}.









\textit{Track record of conference and / or workshop oral and / or poster presentation(s)}

    I have presented my work in a number of conferences and workshops, 
    you may see my cv for the full list.
    A special mention should be made for 
    my talk on Mikrobiokosmos 2022, the yearly conference held by the Greek community for microbiology,
    where I presented my first metabolic modeling related work,
    that also gave me a talk award. 

    I also had the opportunity to present my work at the last two ISME conferences, 
    which greatly deepened my perspective on microbial ecology.









\textit{
    Extensive IT skills, including experience with standard desktop software 
    (e.g. office \& graphics packages), research specific software (e.g. R stats) 
    and online databases (e.g. European Nucleotide Archive, NCBI)
}

 My phd was about data integration, so i can do a thing or two




\textit{
    Proven ability to use initiative to efficiently plan, optimise and 
    progress projects and communicate findings
}


\textit{
    Proven ability to work as part of a collaborative team and 
    to build and maintain effective working relationships with colleagues
}

During my phd, i managed to work with metabolic modeling even if i was not supposed to in the first place. 
I met geomscale , a group of computational geometrists, and together we explored metabolic models and designed 
an algorithm for flux sampling. 

I have developed workflows that support the needs of two european infrustructures lifewatch eric 
(pema: amplicon analysis) and embrc (metagoflow: shotgun metagenome analysis), which i still coordinate 
and support. 

I supervise a set of students not only from the lab that hosts me, but also from other institutions: 
Democritus University of Thrace, Google Summer of Code. 


\textit{Excellent communication and interpersonal skills}


\textit{Excellent time management and organisational skills}



\textit{
    Ability and willingness to contribute to the work of others 
    by offering practical and intellectual help
}

Crete sampling day 


\textit{Ability to work independently and as part of a team}

During my PhD, I was more than happy to co-organize and 



\textit{Ability to present in both written and oral publications}

SITUATION: 
TASK you had to complete: 
ACTION you took:
RESULT: 


\textit{Ability to meet deadlines}



\textit{Ability to contribute to broader management and administrative processes.}


\textit{Ability to assess and organise resources}


\textit{Understand equal opportunity issues as they may impact on areas of research content.}

Equal opportunities are not good enough since we don't have the same starting points. 
It is my belief that much more is needed and i am trying to do something for this through my communities. 

% --------------------------------------------
% \newpage\vspace*{0.3in} 
% --------------------------------------------



\textbf{Desirable Skills, Knowledge and Experience }

\textit{Experience building and analysing metabolic flux models}

As already mentioned, I have developed microbetag~\cite{zafeiropoulos2024microbetag}
and I was among they key developers of dingo~\cite{chalkis2024dingo}.
As mentioned above too, in a recent study of ours, we tested the accuracy of 
GEMs in predicting interaction strengths between pairs of species~\cite{joseph2024predicting}.


Also, I was among the developers of dnngior~\cite{boer2024improving}, a neural network based gapfilling tool. 

Currently, I am working on the analysis of metabolic models of a set of probiotics 
and \textit{Salmonella infantis} to shed some light in the mechanisms with which the former decrease the latter,
if any metabolic interaction is involved.




\textit{Experience with a variety of different machine learning approaches}

In the dnngior~\cite{boer2024improving}, we trained a deep neural network on >11k bacterial species to recover missing reactions.





\textit{
    Experience of directing and supervising practical computer and/or lab-based research 
    of others (e.g. students, technicians)
}

probably the part of the job I enjoy the most. 
Over the last years, I had the chance to supervise a number of students coming from different backgrounds,
and exploring different topics for their bachelor or master thesis.

More specifically, I had the chance to work with \href{https://github.com/ermismd}{Ermis Delopoulos}, 
a master student in the bioinformatics program at KU Leuven.
Ermis had a biological background and no experience in Java at all. 
Me neither. 
Yet, we needed to develop a Java based app to make microbetag accessible on the Cytoscape, 
the most widely used network visualization tool in biology.
In the end, we managed to develop the app,
and Ermis did a great job in his master thesis, which he defended with distinction.


An even more challenging task was the supervision of a set of students over the last three years
with whom we designed and developed the already mentioned {$\mu$}growthDB.
In the first year of my post-doc at KU Leuven, I had the chance to work with
 \href{https://github.com/jcasadogp}{Julia Casado Gómez-Pallete }

then, in the next year 
\href{https://github.com/sofiamonsalve}{Sofia Machado}
\href{https://github.com/AndrewRadev/}{Andrew Radev}



\href{https://github.com/EmmettPeng/}{Xi Peng}



Anna Voukouna 

Last but not least,~\href{https://github.com/SotirisTouliopoulos/}{Sotiris Touliopoulos}
an amazing bachelor student from Democritus University of Thrace, Greece,
he was a member of the IGEM team of his university, and after they contacted me to help them with their project,
he decided to stay and do an after bachelor project with me.
He managed to first enroll as a Google Summer of Code student, in the GeomScale project,
and then continue working with me doing his Erasmus+ traineeship at Prof's Faust lab.
Together, we have been investigating what one can do with the findings of flux sampling, 
and we have developed a Python package, called dingo-stats, for the statistical analysis of flux sampling results.
We are now preparing a manuscript for publication.


% --------------------------------------------


% I am currently a 
% % current position
% postdoctoral researcher at the \href{http://www.msysbiology.com/}{Lab of Microbial Systems Biology}, KU Leuven, 
% working on the~\href{https://www.3domics.eu/}{3D'omics} project, that focuses on poultry and swine gut microbiome. 
% % role
% Specifically, I am involved in the development of microbial network models that incorporate 
% spatial and metabolic data to predict species interactions and evaluate how metabolites 
% influence species distributions.

% My research interests focus on unraveling the \textit{Gordian knot} of microbial ecosystems,
% the laws that govern their structure and dynamics and the role of evolution in shaping them. 
% In the loop of communities shaped by and re-shaping back their habitat, microbial interactions 
% work as the interface where the sum effects of those laws are reflected. 
% Therefore, my ambition is to hack microbial interactions in terms of developing a framework 
% for their interpretation and dynamics, and I am confident that metabolic modeling, 
% when integrated with real-world data, can play a pivotal role in advancing this endeavor.

% % During my PhD, I focused on developing workflows and software tools tailored to 
% % real-world microbial communities, particularly for the analysis of high-throughput 
% % sequencing data. 
% % I also served as a lead developer of a \href{https://prego.hcmr.gr/}{data integration platform} 
% % designed to extract and consolidate associations among species, environments, 
% % and biological processes from both literature and abundance tables from key platforms 
% % such as MGnify, and IMG/JGI.

% Over the last two years, I have been developing a 
% \href{https://microbetag.readthedocs.io/}{microbial co-occurrence network annotator tool}, 
% and deepened my understanding of pairwise metabolic complementarities, 
% particularly within the context of cross-feeding interactions, 
% through topological analysis of their corresponding metabolic models. 
% Further, I have been contributing on the \href{https://geomscale.github.io/}{GeomScale project}
% developing an \href{https://github.com/geomScale/dingo}{efficient framework for flux sampling} 
% within the solution space of metabolic models.
% To support further our metabolic modeling based efforts, I have also been working on 
% the development of a \href{https://mgrowthdb.gbiomed.kuleuven.be/}{microbial growth data repository}; 
% a key resource for constraining models more accurately. 



% % \newpage\vspace*{0.3in} 

% I am confident that I can apply my current expertise to contribute effectively to the project’s goals,
%  while also acquiring new skills in the process—an aspect I find particularly valuable.
% I am enclosing my CV and a letter of recommendation from Prof. Faust. 
% Please do not hesitate to contact me for anything that might be needed. 








